%%%%%%%%%%%%%%%%%%%%%%%%%%%%%%%%%%%%%%%%%%%%%%%%%%%%%%%%%%%%%%%%%%%%%%%%%%%
% Chapter: Active Filters
% Falstad examples: examples/filters (generated by falstad/circuits/filters.py)
%%%%%%%%%%%%%%%%%%%%%%%%%%%%%%%%%%%%%%%%%%%%%%%%%%%%%%%%%%%%%%%%%%%%%%%%%%%

% filfig=<s>: scales coordinates and component size by s, but not the text
\tikzset{filfig/.style={scale=#1, /tikz/circuitikz/bipoles/length=#1cm}}

%%%%%%%%%%%%%%%%%%%%%%%%%%%%%%%%%%%%%%%%%%%%%%%%%%%%%%%%%%%%%%%%%%%%%%%%%%%
\section{Passive RC Filters}
%%%%%%%%%%%%%%%%%%%%%%%%%%%%%%%%%%%%%%%%%%%%%%%%%%%%%%%%%%%%%%%%%%%%%%%%%%%
\nsl{A filter lets some frequencies pass and attenuates others. In a
measurement system the wanted signal is usually slow, while interference
(mains hum at 50\,Hz, switching noise, radio frequencies) is faster. A
low-pass filter in front of the analog-to-digital converter removes it.
It has a second, more fundamental task: a sampling system with sample rate
$f_s$ cannot distinguish a signal at frequency $f$ from one at
$f \pm k f_s$. Everything above $f_s/2$ must be removed \emph{before}
sampling; this is the anti-aliasing filter (chapter~8 explains aliasing in
detail).
\newline
We start with the passive RC filter and the tools to describe it: complex
transfer function, magnitude and phase, decibels and the Bode plot. Then
we add op-amps (chapter~5) to get first-order and second-order active
filters, and finally design the 2.5\,Hz anti-aliasing low-pass of the
capstone task. Further reading: \cite{tietze2008} (chapters on filters),
\cite{horowitz2015} (chapter~6), \cite{mancini2002} (active filter design)
and \cite{sedra2020}.\newpage}

\begin{description}
\item[{\rot\bf RC low-pass}]~\\[-\lblskip]
\begin{center}
\begin{circuitikz}[filfig=1.2]
  \draw (0,0) to[short, o-] (0.5,0) to[R, l=$R$] (3,0) to[short, -o] (4.2,0);
  \draw (3,0) to[C, l=$C$] (3,-2) node[ground]{};
  \draw (0,-2) node[ocirc]{} -- (0,-2);
  \draw (0,0) to[open, v^=$\Vin$] (0,-2);
  \draw (4.2,0) to[open, v^=$\Vout$] (4.2,-2) node[ocirc]{};
  \draw (0,-2) -- (4.2,-2);
\end{circuitikz}
\end{center}
\bi
	\ite low frequency: the capacitor is open, $\Vout = \Vin$
	\ite high frequency: the capacitor is a short, $\Vout \to 0$
	\ite a frequency-dependent voltage divider with the impedance $Z_C = \dfrac{1}{j\omega C}$
\ei
\keyeq{RC low-pass}{$\displaystyle \underline{H}(j\omega) = \frac{\underline{V}_\mathrm{out}}{\underline{V}_\mathrm{in}}
= \frac{1/(j\omega C)}{R + 1/(j\omega C)} = \frac{1}{1 + j\omega RC}$}
\os{\newpage}

\item[{\rot\bf Magnitude, phase and cutoff frequency}]~\\[-\lblskip]
\bi
	\ite with $f_c = \dfrac{1}{2\pi RC}$: \quad $\underline{H} = \dfrac{1}{1 + j f/f_c}$
	\ite magnitude (amplitude ratio) and phase:
\ei
\keyeq{First-order low-pass}{$\displaystyle |\underline{H}| = \frac{1}{\sqrt{1 + (f/f_c)^2}} \qquad
\varphi = -\arctan(f/f_c) \qquad f_c = \frac{1}{2\pi R C}$}
\bi
	\ite at $f = f_c$: $|\underline{H}| = 1/\sqrt{2} = 0.707$, $\varphi = \ang{-45}$
	\ite far above $f_c$: $|\underline{H}| \approx f_c/f$ -- ten times the frequency, one tenth of the amplitude
\ei
\nsl{For a sine input the output of a linear circuit is a sine of the same
frequency. Its amplitude is multiplied by $|\underline{H}|$ and its phase
is shifted by $\varphi$. The cutoff frequency $f_c$ (also corner frequency
or $-3$\,dB frequency) is the frequency at which the real part and the
imaginary part of the denominator are equal; there the output power is half
the power at low frequencies. The time constant $\tau = RC$ and the cutoff
frequency are two views of the same thing: $f_c = 1/(2\pi\tau)$.}
\os{\newpage}

\item[{\rot\bf Decibels}]~\\[-\lblskip]
\bi
	\ite gains and attenuations span many decades: a logarithmic measure is convenient
\ei
\defn{decibel}{The gain in decibels is $\;|H|_\mathrm{dB} = 20\log_{10}|H|$\; (voltage or current ratios).}
\begin{center}
\begin{tabular}{rrcrr}
\toprule
$|H|$ & dB & & $|H|$ & dB \\
\midrule
1 & \SI{0}{\deci\bel} & & 2 & \SI{+6}{\deci\bel} \\
$1/\sqrt{2} = 0.707$ & \SI{-3}{\deci\bel} & & 10 & \SI{+20}{\deci\bel} \\
0.5 & \SI{-6}{\deci\bel} & & 100 & \SI{+40}{\deci\bel} \\
0.1 & \SI{-20}{\deci\bel} & & 1000 & \SI{+60}{\deci\bel} \\
0.01 & \SI{-40}{\deci\bel} & & $10^5$ & \SI{+100}{\deci\bel} \\
\bottomrule
\end{tabular}
\end{center}
\bi
	\ite gains of cascaded stages \textbf{multiply}; in dB they \textbf{add}
	\ite a decade is a factor 10 in frequency, an octave a factor 2
\ei
\os{\newpage}

\item[{\rot\bf Bode plot of the RC low-pass}]~\\[-\lblskip]
\begin{center}
\begin{tikzpicture}
\begin{groupplot}[group style={group size=1 by 2, vertical sep=4mm, xticklabels at=edge bottom},
  width=12.5cm, height=3.9cm, xmode=log, xmin=0.01, xmax=100, grid=major,
  ticklabel style={font=\scriptsize}, label style={font=\small}]
\nextgroupplot[ylabel={$|H|$ in dB}, ymin=-45, ymax=5, ytick={-40,-20,0}]
  \addplot[rwuviolet, very thick, domain=-2:2, samples=100] ({10^x}, {-10*log10(1+10^(2*x))});
  \addplot[rwugray, dashed] coordinates {(0.01,0) (1,0) (100,-40)};
\nextgroupplot[ylabel={$\varphi$ in degrees}, ymin=-95, ymax=5, ytick={-90,-45,0},
  xlabel={$f/f_c$}]
  \addplot[rwucyan, very thick, domain=-2:2, samples=100] ({10^x}, {-atan(10^x)});
\end{groupplot}
\end{tikzpicture}
\end{center}
\bi
	\ite logarithmic frequency axis, magnitude in dB: the curve becomes two straight asymptotes
	\ite \SI{0}{\deci\bel} below $f_c$, \SI{-20}{\deci\bel} per decade above; the corner is at $f_c$ (\SI{-3}{\deci\bel})
\ei
\os{\newpage}

\item[{\rot\bf RC high-pass}]~\\[-\lblskip]
\begin{center}
\begin{circuitikz}[filfig=1.2]
  \draw (0,0) to[short, o-] (0.5,0) to[C, l=$C$] (3,0) to[short, -o] (4.2,0);
  \draw (3,0) to[R, l=$R$] (3,-2) node[ground]{};
  \draw (0,0) to[open, v^=$\Vin$] (0,-2) node[ocirc]{};
  \draw (4.2,0) to[open, v^=$\Vout$] (4.2,-2) node[ocirc]{};
  \draw (0,-2) -- (4.2,-2);
\end{circuitikz}
\end{center}
\keyeq{First-order high-pass}{$\displaystyle \underline{H} = \frac{j f/f_c}{1 + j f/f_c} \qquad
|\underline{H}| = \frac{f/f_c}{\sqrt{1 + (f/f_c)^2}} \qquad \varphi = \ang{90} - \arctan(f/f_c)$}
\bi
	\ite blocks DC (coupling capacitor), \SI{+20}{\deci\bel} per decade below $f_c$
	\ite same cutoff frequency $f_c = 1/(2\pi RC)$, phase $\ang{+45}$ at $f_c$
\ei
\os{\newpage}

\item[{\rot\bf RC filters in Falstad}]~\\[-\lblskip]
\bi
	\ite $R = \SI{1.6}{\kilo\ohm}$, $C = \SI{100}{\nano\farad}$: $f_c = \SI{995}{\hertz}$; three copies driven with \SI{1}{\volt} at \SI{100}{\hertz}, \SI{1}{\kilo\hertz}, \SI{10}{\kilo\hertz}
\ei
\begin{center}
\begin{tabular}{lccc}
\toprule
 & \SI{100}{\hertz} & \SI{1}{\kilo\hertz} & \SI{10}{\kilo\hertz} \\
\midrule
low-pass amplitude & \SI{0.995}{\volt} & \SI{0.705}{\volt} (\ang{-45}) & \SI{0.099}{\volt} \\
high-pass amplitude & \SI{0.100}{\volt} & \SI{0.709}{\volt} (\ang{+45}) & \SI{0.995}{\volt} \\
\bottomrule
\end{tabular}
\end{center}
\falstad{filter_rc_lowpass}\quad\falstad{filter_rc_highpass}\par
\bi
	\ite limits of passive filters: the load changes $f_c$ (it is in parallel to $C$ or $R$), no gain, steep filters need inductors
	\ite remedy: an op-amp buffers the filter and adds gain $\rightarrow$ \textbf{active filter}
\ei
\os{\newpage}
\end{description}

%%%%%%%%%%%%%%%%%%%%%%%%%%%%%%%%%%%%%%%%%%%%%%%%%%%%%%%%%%%%%%%%%%%%%%%%%%%
\section{First-Order Active Filters}
%%%%%%%%%%%%%%%%%%%%%%%%%%%%%%%%%%%%%%%%%%%%%%%%%%%%%%%%%%%%%%%%%%%%%%%%%%%
\nsl{An op-amp turns the RC filter into an active filter: the filter is no
longer loaded by the following stage, and it can have a gain. In the
simplest case the RC network sits in front of a non-inverting amplifier;
alternatively the capacitor is placed into the feedback path of an
inverting amplifier. The cutoff frequency is still $1/(2\pi RC)$ for the
resistor and capacitor that form the filter.\newpage}

\begin{description}
\item[{\rot\bf Non-inverting first-order low-pass}]~\\[-\lblskip]
\begin{center}
\begin{circuitikz}[filfig=1.2]
  \draw (0,0) node[op amp, noinv input up](OA){};
  \coordinate (X) at ($(OA.out)+(0.5,0)$);
  \draw (OA.+) -- ++(-0.5,0) coordinate(c);
  \draw (c) to[R, l_=$R$, -o] ++(-2.2,0) node[left]{$\Vin$};
  \draw (c) to[C, l_=$C$] ++(0,1.4) node[ground, rotate=180]{};
  \draw (OA.-) -- ++(0,-0.9) coordinate(N);
  \draw (N) to[R, l_=$R_2$] (N -| X) -- (X);
  \draw (N) to[R, l_=$R_1$] ++(0,-1.4) node[ground]{};
  \draw (OA.out) -- (X) to[short,-o] ++(0.8,0) node[right]{$\Vout$};
  \draw (c) node[circ]{};
\end{circuitikz}
\end{center}
\bi
	\ite RC low-pass, then a non-inverting amplifier (or a follower for $G=1$)
	\ite $\displaystyle \underline{H} = \left(1+\frac{R_2}{R_1}\right)\frac{1}{1+j\omega RC}$,\quad $f_c = \dfrac{1}{2\pi RC}$
	\ite the op-amp input does not load $C$; the output can drive the next stage
\ei
\os{\newpage}

\item[{\rot\bf Inverting first-order low-pass}]~\\[-\lblskip]
\begin{center}
\begin{circuitikz}[filfig=1.2]
  \draw (0,0) node[op amp](OA){};
  \coordinate (X) at ($(OA.out)+(0.5,0)$);
  \draw (OA.+) -- ++(0,-0.4) node[ground]{};
  \draw (OA.-) -- ++(-0.4,0) coordinate(N);
  \draw (N) to[R, l_=$R_1$, -o] ++(-2.2,0) node[left]{$\Vin$};
  \draw (N) -- ++(0,1.0) coordinate(T) to[R, l_=$R_2$] (T -| X) -- (X);
  \draw (T) -- ++(0,1.1) coordinate(T2) to[C, l=$C$] (T2 -| X) -- (T -| X);
  \draw (OA.out) -- (X) to[short,-o] ++(0.8,0) node[right]{$\Vout$};
  \draw (N) node[circ]{} (T) node[circ]{} (T -| X) node[circ]{};
\end{circuitikz}
\end{center}
\bi
	\ite feedback impedance $Z_2 = R_2 \parallelto \frac{1}{j\omega C}$, gain $-Z_2/R_1$
\ei
\keyeq{Inverting first-order low-pass}{$\displaystyle \underline{H} = -\frac{R_2}{R_1}\,\frac{1}{1 + j\omega R_2 C}
\qquad \text{DC gain } -\frac{R_2}{R_1}, \quad f_c = \frac{1}{2\pi R_2 C}$}
\os{\newpage}

\item[{\rot\bf First-order high-pass (inverting)}]~\\[-\lblskip]
\begin{center}
\begin{circuitikz}[filfig=1.2]
  \draw (0,0) node[op amp](OA){};
  \coordinate (X) at ($(OA.out)+(0.5,0)$);
  \draw (OA.+) -- ++(0,-0.4) node[ground]{};
  \draw (OA.-) -- ++(-0.4,0) coordinate(N);
  \draw (N) to[R, l_=$R_1$] ++(-1.8,0) to[C, l_=$C_1$, -o] ++(-1.6,0) node[left]{$\Vin$};
  \draw (N) -- ++(0,1.0) coordinate(T) to[R, l=$R_2$] (T -| X) -- (X);
  \draw (OA.out) -- (X) to[short,-o] ++(0.8,0) node[right]{$\Vout$};
  \draw (N) node[circ]{};
\end{circuitikz}
\end{center}
\keyeq{Inverting first-order high-pass}{$\displaystyle \underline{H} = -\frac{R_2}{R_1}\,\frac{j\omega R_1 C_1}{1 + j\omega R_1 C_1}
\qquad \text{HF gain } -\frac{R_2}{R_1}, \quad f_c = \frac{1}{2\pi R_1 C_1}$}
\bi
	\ite Falstad: all three circuits with $f_c=\SI{995}{\hertz}$, input \SI{0.1}{\volt} at \SI{1}{\kilo\hertz}
	\itee non-inverting LP ($G=2$): \SI{0.141}{\volt}; inverting LP ($-10$): \SI{0.705}{\volt}; inverting HP ($-10$): \SI{0.709}{\volt}
\ei
\falstad{filter_active_first_order}\par
\os{\newpage}

\item[{\rot\bf Integrator}]~\\[-\lblskip]
\begin{center}
\begin{circuitikz}[filfig=1.2]
  \draw (0,0) node[op amp](OA){};
  \coordinate (X) at ($(OA.out)+(0.5,0)$);
  \draw (OA.+) -- ++(0,-0.4) node[ground]{};
  \draw (OA.-) -- ++(-0.4,0) coordinate(N);
  \draw (N) to[R, l_=$R$, -o] ++(-2.2,0) node[left]{$\Vin$};
  \draw (N) -- ++(0,1.0) coordinate(T) to[C, l_=$C$] (T -| X) -- (X);
  \draw (T) -- ++(0,1.1) coordinate(T2) to[R, l=$R_p$] (T2 -| X) -- (T -| X);
  \draw (OA.out) -- (X) to[short,-o] ++(0.8,0) node[right]{$\Vout$};
  \draw (N) node[circ]{} (T) node[circ]{} (T -| X) node[circ]{};
\end{circuitikz}
\end{center}
\bi
	\ite the current $\Vin/R$ charges $C$: \quad $\displaystyle \Vout(t) = -\frac{1}{RC}\int \Vin\,dt$, \quad $\underline{H} = -\dfrac{1}{j\omega RC}$
	\ite \SI{-20}{\deci\bel} per decade at all frequencies; DC gain $\infty$: any offset drives it into saturation
	\ite practical integrator: $R_p$ limits the DC gain to $-R_p/R$ (it is then a low-pass with $f_c = 1/(2\pi R_p C)$)
\ei
\os{\newpage}

\item[{\rot\bf Integrator: square wave in, triangle out}]~\\[-\lblskip]
\begin{center}
\begin{tikzpicture}
\begin{axis}[width=12cm, height=4.6cm, xmin=0, xmax=4, ymin=-1.3, ymax=1.3,
  xlabel={$t$ in \si{\milli\second}}, ylabel={$V$ in \si{\volt}}, grid=major,
  legend pos=outer north east, legend style={font=\scriptsize},
  ticklabel style={font=\scriptsize}, label style={font=\small}]
  \addplot[rwugray, thick, const plot] coordinates {(0,1) (1,-1) (2,1) (3,-1) (4,-1)};
  \addplot[rwuviolet, very thick] coordinates {(0,0) (1,-1) (2,0) (3,-1) (4,0)};
  \legend{$\Vin$, $\Vout$}
\end{axis}
\end{tikzpicture}
\end{center}
\bi
	\ite $R=\SI{10}{\kilo\ohm}$, $C=\SI{100}{\nano\farad}$, $R_p = \SI{1}{\mega\ohm}$; $\Vin = \SI{\pm 1}{\volt}$, \SI{500}{\hertz}
	\ite slope $\Vin/(RC) = \SI{1000}{\volt\per\second}$: \SI{1}{\volt} in \SI{1}{\milli\second}, triangle \SI{1.0}{\volt} peak-to-peak
	\ite the triangle starts between \SI{0}{\volt} and \SI{-1}{\volt}; $R_p$ centres it slowly ($\tau = R_p C = \SI{0.1}{\second}$)
\ei
\falstad{filter_integrator}\par
\os{\newpage}
\end{description}

%%%%%%%%%%%%%%%%%%%%%%%%%%%%%%%%%%%%%%%%%%%%%%%%%%%%%%%%%%%%%%%%%%%%%%%%%%%
\section{Second-Order Filters}
%%%%%%%%%%%%%%%%%%%%%%%%%%%%%%%%%%%%%%%%%%%%%%%%%%%%%%%%%%%%%%%%%%%%%%%%%%%
\nsl{A first-order filter falls by only 20\,dB per decade. If the
interference is only one decade above the signal, a first-order filter
attenuates it by a factor of ten -- often not enough. Cascading RC stages
helps little: passive RC stages load each other and the corner becomes
soft. Active second-order filters give 40\,dB per decade with one op-amp,
and a sharp corner whose shape is set by the quality factor $Q$. The most
popular circuit is the Sallen-Key filter \cite{mancini2002,tietze2008}.\newpage}

\begin{description}
\item[{\rot\bf Filter order and attenuation}]~\\[-\lblskip]
\begin{center}
\begin{tikzpicture}
\begin{semilogxaxis}[width=12cm, height=5cm, xmin=0.1, xmax=100, ymin=-85, ymax=5,
  xlabel={$f/f_c$}, ylabel={$|H|$ in dB}, grid=major, ytick={-80,-60,-40,-20,0},
  ticklabel style={font=\scriptsize}, label style={font=\small},
  legend style={font=\scriptsize, at={(0.02,0.05)}, anchor=south west}]
  \addplot[rwucyan, very thick, domain=-1:2, samples=100] ({10^x}, {-10*log10(1+10^(2*x))});
  \addplot[rwuviolet, very thick, domain=-1:2, samples=100] ({10^x}, {-10*log10(1+10^(4*x))});
  \addplot[rwucyandark, very thick, dashed, domain=-1:2, samples=100] ({10^x}, {-10*log10(1+10^(8*x))});
  \addplot[rwugray, dashed] coordinates {(20,-85) (20,5)};
  \legend{1st order, 2nd order, 4th order}
\end{semilogxaxis}
\end{tikzpicture}
\end{center}
\bi
	\ite Butterworth filter of order $n$: $|\underline{H}| = 1/\sqrt{1+(f/f_c)^{2n}}$; slope $-20n$\,dB per decade
	\ite dashed line: $f = 20\,f_c$ (\SI{50}{\hertz} for $f_c = \SI{2.5}{\hertz}$): \SI{-26}{\deci\bel}, \SI{-52}{\deci\bel}, \SI{-104}{\deci\bel}
\ei
\os{\newpage}

\item[{\rot\bf Sallen-Key low-pass (unity gain)}]~\\[-\lblskip]
\begin{center}
\begin{circuitikz}[filfig=1.2]
  \draw (0,0) node[op amp, noinv input up](OA){};
  \coordinate (X) at ($(OA.out)+(0.5,0)$);
  \draw (OA.+) -- ++(-0.9,0) coordinate(b);
  \draw (b) to[R, l_=$R_2$] ++(-2,0) coordinate(a) to[R, l_=$R_1$, -o] ++(-2,0) node[left]{$\Vin$};
  \draw (b) to[C, l_=$C_2$] ++(0,-2.0) node[ground]{};
  \draw (a) -- ++(0,1.2) coordinate(t) to[C, l=$C_1$] (t -| X) -- (X);
  \draw (OA.-) -- ++(0,-0.6) -| ($(OA.out)+(0.25,0)$);
  \draw (OA.out) -- (X) to[short,-o] ++(0.8,0) node[right]{$\Vout$};
  \draw (a) node[circ]{} (b) node[circ]{} (X) node[circ]{};
\end{circuitikz}
\end{center}
\bi
	\ite the follower makes $\Vout = V_+$; $C_1$ feeds the output back to the middle node (positive feedback near $f_0$)
	\ite $\displaystyle \underline{H} = \frac{1}{1 + j\omega C_2(R_1+R_2) - \omega^2 R_1R_2C_1C_2}$
\ei
\os{\newpage}

\item[{\rot\bf Natural frequency and quality factor}]~\\[-\lblskip]
\bi
	\ite comparing with $1/(1 + j\frac{f}{Q f_0} - (f/f_0)^2)$ gives $f_0$ and $Q$:
\ei
\keyeq{Unity-gain Sallen-Key low-pass}{$\displaystyle f_0 = \frac{1}{2\pi\sqrt{R_1R_2C_1C_2}} \qquad
Q = \frac{\sqrt{R_1R_2C_1C_2}}{C_2\,(R_1+R_2)} \quad\stackrel{R_1=R_2}{=}\quad \frac12\sqrt{\frac{C_1}{C_2}}$}
\begin{center}
\begin{tikzpicture}
\begin{semilogxaxis}[width=11cm, height=4.4cm, xmin=0.1, xmax=10, ymin=-40, ymax=8,
  xlabel={$f/f_0$}, ylabel={$|H|$ in dB}, grid=major, ytick={-40,-20,0},
  ticklabel style={font=\scriptsize}, label style={font=\small},
  legend style={font=\scriptsize, at={(0.02,0.05)}, anchor=south west}]
  \addplot[rwucyan, very thick, domain=-1:1, samples=120] ({10^x}, {-10*log10((1-10^(2*x))^2+10^(2*x)/0.25)});
  \addplot[rwuviolet, very thick, domain=-1:1, samples=120] ({10^x}, {-10*log10((1-10^(2*x))^2+10^(2*x)/0.5)});
  \addplot[rwucyandark, very thick, dashed, domain=-1:1, samples=120] ({10^x}, {-10*log10((1-10^(2*x))^2+10^(2*x)/4)});
  \legend{$Q=0.5$, $Q=0.707$ (Butterworth), $Q=2$}
\end{semilogxaxis}
\end{tikzpicture}
\end{center}
\bi
	\ite $|\underline{H}(f_0)| = Q$; Butterworth $Q = 1/\sqrt{2}$: maximally flat, \SI{-3}{\deci\bel} at $f_0$; higher $Q$: peaking and ringing
\ei
\os{\newpage}

\item[{\rot\bf Design of a Butterworth Sallen-Key low-pass}]~\\[-\lblskip]
\bi
	\ite capacitors come in coarse steps: choose them first, then calculate the resistors
\ei
\designcard{Unity-gain Sallen-Key low-pass, $Q = 0.707$}{
\begin{enumerate}
\item Choose $R_1 = R_2 = R$. Then $Q = 0.707$ requires $C_1 = 2\,C_2$ ($C_1$ goes to the output).
\item Choose $C_2$ from the E6 series, $C_1 \approx 2 C_2$ (e.g.\ \SI{10}{\nano\farad} and \SI{22}{\nano\farad}: $Q = 0.74$).
\item Calculate $R = \dfrac{1}{2\pi f_0 \sqrt{C_1 C_2}}$ and round to the E24 or E96 series.
\item Check $f_0$ and $Q$ with the rounded values. Use film or C0G capacitors (stable, low leakage).
\end{enumerate}}
\bi
	\ite example $f_0 = \SI{1}{\kilo\hertz}$: $C_2 = \SI{10}{\nano\farad}$, $C_1 = \SI{22}{\nano\farad}$, $R = \SI{10.7}{\kilo\ohm}$ (E96): $f_0 = \SI{1003}{\hertz}$, $Q = 0.742$
\ei
\nsl{Capacitors are available in coarser steps than resistors, so they are
chosen first and the resistors are calculated. Exact Butterworth
behaviour is rarely needed; $Q$ between 0.65 and 0.75 changes the response
near $f_0$ by less than 1\,dB. Higher-order filters are built as a cascade
of second-order stages (and one first-order stage for odd orders) with
tabulated values of $f_0$ and $Q$ for each stage \cite{tietze2008}.}
\os{\newpage}

\item[{\rot\bf First order vs.\ second order in Falstad}]~\\[-\lblskip]
\begin{center}
\begin{tabular}{lcc}
\toprule
 & \SI{1}{\kilo\hertz} & \SI{10}{\kilo\hertz} \\
\midrule
first order RC ($f_c = \SI{995}{\hertz}$) & \SI{0.705}{\volt} & \SI{0.099}{\volt} (\SI{-20}{\deci\bel}) \\
Sallen-Key ($f_0 = \SI{1003}{\hertz}$, $Q=0.742$) & \SI{0.744}{\volt} & \SI{0.0101}{\volt} (\SI{-40}{\deci\bel}) \\
\bottomrule
\end{tabular}
\end{center}
\bi
	\ite input \SI{1}{\volt} amplitude; one decade above the corner the second-order filter attenuates ten times more
	\ite near $f_0$ both behave similarly
\ei
\falstad{filter_sallen_key}\par
\os{\newpage}
\end{description}

%%%%%%%%%%%%%%%%%%%%%%%%%%%%%%%%%%%%%%%%%%%%%%%%%%%%%%%%%%%%%%%%%%%%%%%%%%%
\section{The Anti-Aliasing Filter of the Capstone}\label{sec:filters-aa}
%%%%%%%%%%%%%%%%%%%%%%%%%%%%%%%%%%%%%%%%%%%%%%%%%%%%%%%%%%%%%%%%%%%%%%%%%%%
\nsl{In the capstone the amplified bridge signal (0 to 4.00\,V, chapter~5)
is sampled by a sample-and-hold circuit with $f_s = 5\,$Hz and converted by
a 3-bit flash ADC. The sensor signal itself changes slowly. The problem is
the 50\,Hz interference of 5\,mV in series with the Pt100: it is a
differential signal (4.9\,mV at the bridge tap), so the instrumentation
amplifier amplifies it by the full gain of 107.6 to an amplitude of
0.52\,V. This section designs the low-pass filter that removes it, with
the values used in the capstone chapter: $R = \SI{6.8}{\kilo\ohm}$,
$C = \SI{10}{\micro\farad}$.\newpage}

\begin{description}
\item[{\rot\bf Why filter before sampling? (preview of chapter 8)}]~\\[-\lblskip]
\bi
	\ite sampling theorem (Nyquist, Shannon \cite{shannon1949}): only frequencies below $f_s/2$ are reproduced correctly
	\ite with $f_s = \SI{5}{\hertz}$: $f_s/2 = \SI{2.5}{\hertz}$ $\Rightarrow$ cutoff of the low-pass $f_c \approx \SI{2.5}{\hertz}$
	\ite a component above $f_s/2$ is not removed by sampling -- it appears at a wrong, low frequency (\textbf{aliasing})
	\itee \SI{50}{\hertz} $= 10\cdot f_s$: the samples see the interference always at the same phase -- a constant, random-looking error
	\itee if the mains frequency drifts to \SI{50.1}{\hertz}, the error becomes a slow \SI{0.1}{\hertz} wave in the data
	\ite after sampling nothing can separate the alias from the signal: the filter must act \textbf{before} the S\&H
\ei
\os{\newpage}

\item[{\rot\bf How much attenuation do we get at 50\,Hz?}]~\\[-\lblskip]
\bi
	\ite first order, $f_c = \SI{2.5}{\hertz}$: $|\underline{H}(\SI{50}{\hertz})| = \dfrac{1}{\sqrt{1+(50/2.5)^2}} = \dfrac{1}{\sqrt{401}} = 0.050$ $\;\hat{=}\;$ \SI{-26}{\deci\bel}
	\ite second order (Butterworth): $\dfrac{1}{\sqrt{1+(50/2.5)^4}} = 0.0025$ $\;\hat{=}\;$ \SI{-52}{\deci\bel}
	\ite capstone filter \SI{6.8}{\kilo\ohm}, \SI{10}{\micro\farad}: $f_c = \SI{2.34}{\hertz}$, $|\underline{H}| = 1/\sqrt{1+(50/2.34)^2} = 0.047$ $\;\hat{=}\;$ \SI{-26.6}{\deci\bel}
	\ite interference before the filter: $\SI{4.875}{\milli\volt}\cdot 107.6 = \SI{0.52}{\volt}$
\ei
\begin{center}
\begin{tabular}{lccc}
\toprule
 & amplitude & in LSB (\SI{0.5}{\volt}) & \\
\midrule
no filter & \SI{0.52}{\volt} & 1 LSB & useless \\
first order, \SI{2.34}{\hertz} & \SI{25}{\milli\volt} & 0.05 LSB & sufficient \\
Sallen-Key, \SI{2.55}{\hertz} & \SI{1.4}{\milli\volt} & 0.003 LSB & large margin \\
\bottomrule
\end{tabular}
\end{center}
\nsl{For a 3-bit converter with 4\,V full scale one LSB is 0.5\,V, and the
quantization error is $\pm 0.25$\,V. A residual of 25\,mV is far below
$\LSB/2$, so the first-order filter of the capstone is sufficient. A
second-order filter costs only one more capacitor and gives a large margin,
for example for a higher-resolution ADC: with 10 bits (LSB = 3.9\,mV) the
first-order filter leaves about 6 LSB of interference, the second-order
filter less than half an LSB.}
\os{\newpage}

\item[{\rot\bf First-order anti-aliasing filter with E-series parts}]~\\[-\lblskip]
\begin{center}
\begin{circuitikz}[filfig=1.1]
  \draw (0,0) node[op amp, noinv input up](OA){};
  \draw (OA.+) -- ++(-0.5,0) coordinate(c);
  \draw (c) to[R, l_=\SI{6.8}{\kilo\ohm}, -o] ++(-2.4,0) node[left]{from INA};
  \draw (c) to[C, l_=\SI{10}{\micro\farad}] ++(0,1.4) node[ground, rotate=180]{};
  \draw (OA.-) -- ++(0,-0.6) -| ($(OA.out)+(0.3,0)$);
  \draw (OA.out) to[short,-o] ++(1,0) node[right]{to S\&H};
  \draw (c) node[circ]{};
\end{circuitikz}
\end{center}
\bi
	\ite $RC = 1/(2\pi\cdot\SI{2.5}{\hertz}) = \SI{63.7}{\milli\second}$; choose $C = \SI{10}{\micro\farad}$ $\Rightarrow$ $R = \SI{6.37}{\kilo\ohm}$
	\ite round so that $f_c \le f_s/2$: \SI{6.8}{\kilo\ohm} (E12) $\Rightarrow$ $f_c = \SI{2.34}{\hertz}$, \SI{-26.6}{\deci\bel} at \SI{50}{\hertz}
	\ite also usable: \SI{62}{\kilo\ohm} with \SI{1}{\micro\farad} ($f_c = \SI{2.57}{\hertz}$, slightly above $f_s/2$, still \SI{-26}{\deci\bel} at \SI{50}{\hertz})
	\ite the follower decouples the filter from the S\&H input
	\ite step response: $\tau = \SI{68}{\milli\second}$; settling to \SI{1}{\percent} takes $4.6\,\tau = \SI{0.31}{\second}$
\ei
\os{\newpage}

\item[{\rot\bf Second-order version (Sallen-Key)}]~\\[-\lblskip]
\begin{center}
\begin{circuitikz}[filfig=1.1]
  \draw (0,0) node[op amp, noinv input up](OA){};
  \coordinate (X) at ($(OA.out)+(0.5,0)$);
  \draw (OA.+) -- ++(-0.9,0) coordinate(b);
  \draw (b) to[R, l_=\SI{91}{\kilo\ohm}] ++(-2,0) coordinate(a) to[R, l_=\SI{91}{\kilo\ohm}, -o] ++(-2,0) node[left]{from INA};
  \draw (b) to[C, l_=\SI{470}{\nano\farad}] ++(0,-2.0) node[ground]{};
  \draw (a) -- ++(0,1.2) coordinate(t) to[C, l=\SI{1}{\micro\farad}] (t -| X) -- (X);
  \draw (OA.-) -- ++(0,-0.6) -| ($(OA.out)+(0.25,0)$);
  \draw (OA.out) -- (X) to[short,-o] ++(0.8,0) node[right]{to S\&H};
  \draw (a) node[circ]{} (b) node[circ]{} (X) node[circ]{};
\end{circuitikz}
\end{center}
\bi
	\ite $C_1 = \SI{1}{\micro\farad}$, $C_2 = \SI{470}{\nano\farad}$: $Q = \frac12\sqrt{1/0.47} = 0.729$
	\ite $R = 1/(2\pi\cdot\SI{2.5}{\hertz}\cdot\sqrt{\SI{1}{\micro\farad}\cdot\SI{470}{\nano\farad}}) = \SI{92.8}{\kilo\ohm}$ $\rightarrow$ E24 \SI{91}{\kilo\ohm}: $f_0 = \SI{2.55}{\hertz}$
\ei
\os{\newpage}

\item[{\rot\bf Both filters in Falstad}]~\\[-\lblskip]
\begin{center}
\begin{tikzpicture}
\begin{axis}[width=12cm, height=4.4cm, xmin=0, xmax=0.6, ymin=1.3, ymax=2.7,
  xlabel={$t$ in \si{\second}}, ylabel={$V$ in \si{\volt}}, grid=major,
  legend pos=outer north east, legend style={font=\scriptsize},
  ticklabel style={font=\scriptsize}, label style={font=\small}]
  \addplot[rwugray, domain=0:0.6, samples=600] {2+sin(360*50*x)};
  \addplot[rwuviolet, very thick, domain=0:0.6, samples=300] {2*(1-exp(-x/0.068))+0.0245*sin(360*50*x-87)};
  \legend{$\Vin$, 1st order}
\end{axis}
\end{tikzpicture}
\end{center}
\bi
	\ite input: \SI{2}{\volt} signal + \SI{0.52}{\volt}, \SI{50}{\hertz} interference (= \SI{4.875}{\milli\volt} $\times$ 107.6)
	\ite after settling: both outputs \SI{2.00}{\volt} DC; remaining \SI{50}{\hertz}: \SI{25}{\milli\volt} (1st order), \SI{1.4}{\milli\volt} (Sallen-Key)
\ei
\falstad{filter_antialias}\par
\os{\newpage}

\item[{\rot\bf Alternative: twin-T notch at 50\,Hz}]~\\[-\lblskip]
\begin{center}
\begin{circuitikz}[filfig=1.0]
  \draw (0,0) to[short, o-] (0.6,0) coordinate(i);
  \draw (i) -- ++(0,1.4) to[R, l=$R$] ++(2.4,0) coordinate(m1) to[R, l=$R$] ++(2.4,0) -- ++(0,-1.4) coordinate(o);
  \draw (i) -- ++(0,-1.0) to[C, l=$C$] ++(2.4,0) coordinate(m2) to[C, l=$C$] ++(2.4,0) -- ++(0,1.0);
  \draw (m1) to[C, l=$2C$] ++(0,-1.1) node[ground]{};
  \draw (m2) to[R, l=$R/2$] ++(0,-1.3) node[ground]{};
  \draw (o) to[short, -o] ++(1,0) node[right]{$\Vout$ (to a follower)};
  \node[left] at (0,0) {$\Vin$};
  \draw (i) node[circ]{} (o) node[circ]{};
\end{circuitikz}
\end{center}
\bi
	\ite zero transmission at $f_0 = 1/(2\pi RC)$: $R = \SI{31.8}{\kilo\ohm}$, $C = \SI{100}{\nano\farad}$ $\Rightarrow$ \SI{50}{\hertz}
	\ite Falstad: \SI{50}{\hertz} reduced from \SI{1}{\volt} to \SI{0.45}{\milli\volt}; but \SI{5}{\hertz} is also attenuated (0.927)
	\ite needs precise parts, removes only one frequency -- it does \textbf{not} replace the anti-aliasing low-pass
\ei
\falstad{filter_twin_t}\par
\os{\newpage}

\item[{\rot\bf Design card}]~\\[-\lblskip]
\designcard{Anti-aliasing low-pass}{
\begin{enumerate}
\item Cutoff frequency from the sample rate: $f_c \le f_s/2$ (capstone: $f_s = 5\,\mathrm{Hz}$, $f_c = 2.34\,\mathrm{Hz}$).
\item Required attenuation: interference amplitude at the ADC input compared with $\LSB/2$; with $f/f_c = 20$ a first-order filter gives $-26\,\mathrm{dB}$, a second-order filter $-52\,\mathrm{dB}$ (capstone: $0.52\,\mathrm{V}\cdot 0.047 = 25\,\mathrm{mV} < 250\,\mathrm{mV}$).
\item First order: $RC = 1/(2\pi f_c)$; choose $C$, calculate $R$, round to the E-series so that $f_c \le f_s/2$ (capstone: $6.8\,\mathrm{k\Omega}$, $10\,\mu\mathrm{F}$). Buffer the output.
\item Second order: unity-gain Sallen-Key with $C_1 \approx 2C_2$, $R = 1/(2\pi f_0\sqrt{C_1C_2})$.
\item Check: settling time of a step ($\approx 5\tau$) against the sample period, and headroom of the amplifier before the filter (signal + interference must not clip).
\end{enumerate}}
\os{\newpage}
\end{description}

%%%%%%%%%%%%%%%%%%%%%%%%%%%%%%%%%%%%%%%%%%%%%%%%%%%%%%%%%%%%%%%%%%%%%%%%%%%
\section{Exercises}
%%%%%%%%%%%%%%%%%%%%%%%%%%%%%%%%%%%%%%%%%%%%%%%%%%%%%%%%%%%%%%%%%%%%%%%%%%%

\exercise{RC low-pass}\label{ex:filt-rc}
An RC low-pass has $R = \SI{1.6}{\kilo\ohm}$ and $C = \SI{100}{\nano\farad}$.
\begin{talist}
\ta Calculate the cutoff frequency.
\ta Calculate magnitude (also in dB) and phase at \SI{100}{\hertz}, \SI{1}{\kilo\hertz} and \SI{10}{\kilo\hertz}.
\ta What changes if a load of \SI{1.6}{\kilo\ohm} is connected to the output?
\ta Check b) in Falstad.
\end{talist}

\begin{solution}
\begin{talist}
\ta $f_c = 1/(2\pi\cdot\SI{1.6}{\kilo\ohm}\cdot\SI{100}{\nano\farad}) = \SI{995}{\hertz}$.
\ta \SI{100}{\hertz}: $f/f_c = 0.1005$, $|H| = 1/\sqrt{1.0101} = 0.995$ (\SI{-0.04}{\deci\bel}), $\varphi = \ang{-5.7}$.\newline
    \SI{1}{\kilo\hertz}: $f/f_c = 1.005$, $|H| = 0.705$ (\SI{-3.0}{\deci\bel}), $\varphi = \ang{-45.2}$.\newline
    \SI{10}{\kilo\hertz}: $f/f_c = 10.05$, $|H| = 0.099$ (\SI{-20.1}{\deci\bel}), $\varphi = \ang{-84.3}$.
\ta The load is in parallel to $C$. Thevenin: $V_\mathrm{th} = \Vin/2$, $R_\mathrm{th} = \SI{0.8}{\kilo\ohm}$:
    the DC gain drops to 0.5 and $f_c$ doubles to \SI{1.99}{\kilo\hertz}. A follower avoids this.
\ta \falstad{filter_rc_lowpass}\newline amplitudes \SI{0.995}{\volt}, \SI{0.705}{\volt} (lagging by \ang{45}, i.e.\ \SI{125}{\micro\second}) and \SI{0.099}{\volt}.
\end{talist}
\end{solution}

\exercise{First-order active filters}\label{ex:filt-active}
\begin{talist}
\ta Design an inverting low-pass with DC gain $-10$, input resistance \SI{10}{\kilo\ohm} and $f_c = \SI{1}{\kilo\hertz}$.
\ta Calculate the output amplitude for \SI{0.1}{\volt} at \SI{1}{\kilo\hertz} with $C$ rounded to the E-series.
\ta Design a non-inverting low-pass with $G=2$ and the same $f_c$, using $C = \SI{10}{\nano\farad}$.
\ta Design an inverting high-pass with HF gain $-10$, $R_1 = \SI{10}{\kilo\ohm}$, $f_c\approx\SI{1}{\kilo\hertz}$.
\ta Check b) to d) in Falstad (input \SI{0.1}{\volt}, \SI{1}{\kilo\hertz}).
\end{talist}

\begin{solution}
\begin{talist}
\ta $R_1 = \SI{10}{\kilo\ohm}$, $R_2 = \SI{100}{\kilo\ohm}$, $C = 1/(2\pi\cdot\SI{100}{\kilo\ohm}\cdot\SI{1}{\kilo\hertz}) = \SI{1.59}{\nano\farad}$.
\ta $C = \SI{1.6}{\nano\farad}$ (E24): $f_c = \SI{995}{\hertz}$; $\SI{0.1}{\volt}\cdot 10\cdot 0.705 = \SI{0.705}{\volt}$.
\ta $R = 1/(2\pi\cdot\SI{1}{\kilo\hertz}\cdot\SI{10}{\nano\farad}) = \SI{15.9}{\kilo\ohm}$ $\rightarrow$ \SI{16}{\kilo\ohm}; gain stage $R_1=R_2=\SI{10}{\kilo\ohm}$.
    Output: $\SI{0.1}{\volt}\cdot 2\cdot 0.705 = \SI{0.141}{\volt}$.
\ta $R_2 = \SI{100}{\kilo\ohm}$, $C_1 = 1/(2\pi\cdot\SI{10}{\kilo\ohm}\cdot\SI{1}{\kilo\hertz}) = \SI{15.9}{\nano\farad}$ $\rightarrow$ \SI{16}{\nano\farad}:
    $f_c = \SI{995}{\hertz}$, output $\SI{0.1}{\volt}\cdot 10\cdot 0.709 = \SI{0.709}{\volt}$.
\ta \falstad{filter_active_first_order}\newline \texttt{Vout\_a} $= \SI{0.141}{\volt}$, \texttt{Vout\_b} $= \SI{0.705}{\volt}$, \texttt{Vout\_c} $= \SI{0.709}{\volt}$.
\end{talist}
\end{solution}

\exercise{Integrator}\label{ex:filt-int}
An integrator has $R = \SI{10}{\kilo\ohm}$, $C = \SI{100}{\nano\farad}$ and $R_p = \SI{1}{\mega\ohm}$ parallel to $C$.
The input is a square wave of \SI{\pm 1}{\volt} and \SI{500}{\hertz}, starting with $+\SI{1}{\volt}$ at $t=0$, $\Vout(0) = 0$.
\begin{talist}
\ta Calculate the slope of the output voltage and sketch $\Vout(t)$ for the first two periods.
\ta Calculate the peak-to-peak value of the triangle.
\ta What is the purpose of $R_p$? What happens without it if the op-amp has an offset of \SI{1}{\milli\volt}?
\ta Above which frequency does the circuit behave as an integrator?
\ta Check a) and b) in Falstad.
\end{talist}

\begin{solution}
\begin{talist}
\ta $d\Vout/dt = -\Vin/(RC) = -\SI{1}{\volt}/\SI{1}{\milli\second} = \SI{-1000}{\volt\per\second}$ while $\Vin = +\SI{1}{\volt}$.
    Triangle: $0 \rightarrow \SI{-1}{\volt}$ in the first \SI{1}{\milli\second}, back to \SI{0}{\volt} at \SI{2}{\milli\second}, and so on.
\ta \SI{1}{\volt} per half period: \SI{1.0}{\volt} peak-to-peak.
\ta $R_p$ limits the DC gain to $-R_p/R = -100$. Without it, \SI{1}{\milli\volt} offset causes a current of \SI{0.1}{\micro\ampere}
    that charges $C$ at \SI{1}{\volt\per\second}: the output saturates after some seconds. With $R_p$: DC error \SI{0.1}{\volt}.
\ta Above $f_c = 1/(2\pi R_p C) = \SI{1.6}{\hertz}$; at \SI{500}{\hertz} the error is negligible.
\ta \falstad{filter_integrator}\newline the output reaches about \SI{-1}{\volt} after \SI{1}{\milli\second}
    (\SI{-0.98}{\volt} in the simulation, $R_p$ and the time step cost a little), the triangle has \SI{1.0}{\volt} peak-to-peak and drifts slowly towards a symmetric position.
\end{talist}
\end{solution}

\exercise{Sallen-Key low-pass}\label{ex:filt-sk}
Design a unity-gain Sallen-Key low-pass with $f_0 = \SI{1}{\kilo\hertz}$ and Butterworth characteristic.
\begin{talist}
\ta Start with $C_2 = \SI{10}{\nano\farad}$ and $R_1 = R_2$. Choose $C_1$ from the E6 series and calculate $R$ (E96).
\ta Calculate $Q$ and $f_0$ with the chosen values.
\ta Calculate the amplitude for \SI{1}{\volt} input at \SI{1}{\kilo\hertz} and \SI{10}{\kilo\hertz}. Compare with a first-order filter.
\ta Check c) in Falstad.
\end{talist}

\begin{solution}
\begin{talist}
\ta $Q = 0.707$ needs $C_1 = 2C_2 = \SI{20}{\nano\farad}$; E6: \SI{22}{\nano\farad}.
    \newline $R = 1/(2\pi\cdot\SI{1}{\kilo\hertz}\cdot\sqrt{\SI{22}{\nano\farad}\cdot\SI{10}{\nano\farad}}) = \SI{10.73}{\kilo\ohm}$ $\rightarrow$ \SI{10.7}{\kilo\ohm}.
\ta $Q = \frac12\sqrt{22/10} = 0.742$; $f_0 = 1/(2\pi\cdot\SI{10.7}{\kilo\ohm}\cdot\SI{14.83}{\nano\farad}) = \SI{1003}{\hertz}$.
\ta $|H| = 1/\sqrt{(1-x^2)^2 + (x/Q)^2}$ with $x = f/f_0$.
    \SI{1}{\kilo\hertz}: $x = 0.997$, $|H| = 0.744$ (\SI{0.744}{\volt}; first order: \SI{0.705}{\volt}).
    \SI{10}{\kilo\hertz}: $x = 9.97$, $|H| = 1/\sqrt{98.4^2 + 13.4^2} = 0.0101$ (\SI{-40}{\deci\bel}; first order: 0.099, \SI{-20}{\deci\bel}).
\ta \falstad{filter_sallen_key}\newline \texttt{VSK\_1kHz} $= \SI{0.744}{\volt}$, \texttt{VSK\_10kHz} $= \SI{10.1}{\milli\volt}$,\newline
    \texttt{V1st\_1kHz} $= \SI{0.705}{\volt}$, \texttt{V1st\_10kHz} $= \SI{0.099}{\volt}$.
\end{talist}
\end{solution}

\exercise{Capstone: anti-aliasing low-pass}\label{ex:filt-aa}
The capstone samples with $f_s = \SI{5}{\hertz}$. The amplifier has a gain of 107.6
and its output carries a \SI{50}{\hertz} interference (\SI{4.875}{\milli\volt} at the bridge tap).
The ADC has 3 bits and \SI{4}{\volt} full scale.
\begin{talist}
\ta Choose the cutoff frequency of the anti-aliasing filter and justify it.
\ta Design a first-order filter (follower behind an RC) with $C = \SI{10}{\micro\farad}$ and an E12 resistor.
\ta Calculate the attenuation at \SI{50}{\hertz} in dB for an ideal first-order and second-order filter with $f_c = \SI{2.5}{\hertz}$, and for your filter.
\ta Calculate the remaining interference at the ADC input with your filter and compare it with $\LSB/2$.
\ta How long does the filter output need to settle within \SI{1}{\percent} after a step of the sensor signal?
\ta Check the result in Falstad (signal \SI{2}{\volt} DC with \SI{0.52}{\volt} at \SI{50}{\hertz}).
\end{talist}

\begin{solution}
\begin{talist}
\ta Nyquist: only frequencies below $f_s/2 = \SI{2.5}{\hertz}$ can be sampled without aliasing, so $f_c \le \SI{2.5}{\hertz}$.
    The temperature changes much more slowly than that.
\ta $R = 1/(2\pi\cdot\SI{2.5}{\hertz}\cdot\SI{10}{\micro\farad}) = \SI{6.37}{\kilo\ohm}$ $\rightarrow$ \SI{6.8}{\kilo\ohm}
    (rounded up, so that $f_c$ stays below $f_s/2$): $f_c = 1/(2\pi\cdot\SI{6.8}{\kilo\ohm}\cdot\SI{10}{\micro\farad}) = \SI{2.34}{\hertz}$.
    (\SI{62}{\kilo\ohm} with \SI{1}{\micro\farad}, $f_c = \SI{2.57}{\hertz}$, slightly above $f_s/2$, would also be acceptable.)
\ta Ideal first order: $1/\sqrt{1+20^2} = 0.0499$, \SI{-26}{\deci\bel}; second order: $1/\sqrt{1+20^4} = 0.0025$, \SI{-52}{\deci\bel}.
    With $f_c = \SI{2.34}{\hertz}$: $1/\sqrt{1+(50/2.34)^2} = 0.047$, \SI{-26.6}{\deci\bel}.
\os{\newpage}
\ta Before the filter: $\SI{4.875}{\milli\volt}\cdot 107.6 = \SI{0.52}{\volt}$; after it: $\SI{0.52}{\volt}\cdot 0.047 = \SI{25}{\milli\volt}$
    $< \LSB/2 = \SI{250}{\milli\volt}$ (0.05 LSB). Sufficient.
\ta $\tau = RC = \SI{68}{\milli\second}$; \SI{1}{\percent}: $t = \tau\ln 100 = 4.6\,\tau = \SI{0.31}{\second}$, i.e.\ about 1.5 sample periods.
\ta \falstad{filter_antialias}\newline after about \SI{0.6}{\second} \texttt{Vout\_1st} shows \SI{2.00}{\volt} with
    \SI{25}{\milli\volt} of \SI{50}{\hertz} ripple; the Sallen-Key output (\SI{91}{\kilo\ohm}, \SI{1}{\micro\farad}, \SI{470}{\nano\farad})
    leaves only \SI{1.4}{\milli\volt}.
\end{talist}
\end{solution}

\exercise{Twin-T notch}\label{ex:filt-twint}
\begin{talist}
\ta Dimension a twin-T notch for \SI{50}{\hertz} with $C = \SI{100}{\nano\farad}$.
\ta Its transfer function is $\underline{H} = \dfrac{1-x^2}{1-x^2+j4x}$ with $x = f/f_0$. Calculate $|H|$ at \SI{5}{\hertz} and at \SI{50.5}{\hertz}.
\ta Compare with the first-order low-pass with $f_c = \SI{2.34}{\hertz}$. Which filter would you use in the capstone?
\ta Check a) and b) in Falstad.
\end{talist}

\begin{solution}
\begin{talist}
\ta $R = 1/(2\pi\cdot\SI{50}{\hertz}\cdot\SI{100}{\nano\farad}) = \SI{31.8}{\kilo\ohm}$, $R/2 = \SI{15.9}{\kilo\ohm}$, $2C = \SI{200}{\nano\farad}$.
\ta \SI{5}{\hertz}: $x = 0.1$, $|H| = 0.99/\sqrt{0.99^2 + 0.4^2} = 0.927$.
    \newline \SI{50.5}{\hertz}: $x = 1.01$, $|H| = 0.0201/\sqrt{0.0201^2+4.04^2} = 0.0050$.
\ta At \SI{50}{\hertz} the notch is much better (\SI{-46}{\deci\bel} at \SI{50.5}{\hertz} and more at exactly \SI{50}{\hertz}),
    but it attenuates the signal at \SI{5}{\hertz} as well, it needs precise components, and it does nothing against
    other frequencies above $f_s/2$. The anti-aliasing low-pass is required anyway; the notch is at most an addition.
\ta \falstad{filter_twin_t}\newline \texttt{Vout\_5Hz} $= \SI{0.927}{\volt}$, \texttt{Vout\_50Hz} $= \SI{0.45}{\milli\volt}$ (input \SI{1}{\volt}).
\end{talist}
\end{solution}

%%% Local Variables:
%%% mode: latex
%%% TeX-master: "docu"
%%% End:
